\documentclass[10pt]{article}

\usepackage[utf8]{inputenc}

\usepackage[T1]{fontenc}

\usepackage{amsmath}

\usepackage{amsfonts}

\usepackage{amssymb}

\usepackage[version=4]{mhchem}

\usepackage{stmaryrd}



\begin{document}

к р и т и ч е с им для функционала $F$, если $\partial F / \partial \gamma=0$ для любой вариации $\partial \gamma$ пути $\gamma_{0}$. При опрсделенных предположениях о функционале $F$ можно определить и нде к с $\lambda$ критич. пути $\gamma_{0}$, опять копируя «конечномерное» определение и беря в качестве $\lambda$ размерность максимального подпространства в пространстве $T_{\gamma} \Omega$, на к-ром гессиан $\partial^{2} F$ отрицательно определен. Таким образом, «критич. точки» $\gamma$ для функционала $F$ находятся как решения дифференциального уравнения $\partial F / \partial \gamma=0$, называемого у р а н е н и ем Эй ле р а - Л а гр а н жа. Среди важнейших функционалов $F$, естественно возникающих в механике, математич. физике, следует отметить: действие пути $\gamma$, то есть



\[

E(\gamma)=\int_{0}^{1}|d \gamma / d t|^{2} d t

\]



и длина пути $\gamma$, то есть



\[

L(\gamma)=\int_{0}^{1}|d \gamma / d t| d t

\]



Т е о рема. Путь $\gamma_{0}$ является критическим («критич. точкой») для функционала $E(\gamma)$ тогда и только тогда, когда $\gamma_{0}-$ это геодезическая (при этом параметром $t$ вдоль геодезической служит длина дуги).



Так как при определении функционалов $E$ и $L$ предполагается наличие на многообразии римановой метрики, то геодезические могут быть также охарактеризованы как кривые локально минимальной длины. Экс т р е ма л и (к р и тиче ские пут и) для функционала длины $L$ - это пути, получающиеся из геодезических $\gamma(t)$ гладкими заменами параметра $t$. Векторное поле $v(t)$ вдоль геодезической $\gamma_{0}$ называется я к о б и вым, если оно удовлетворяет д и фференциальному уравнению Якоби:



\[

\left(\nabla_{\dot{\gamma}_{0}}\right)^{2} v+R\left(\dot{\gamma}_{0}, v\right) \dot{\gamma}_{0}=0

\]



где $\nabla$ - оператор ковариантного дифференцирования, $R-$ оператор, соответствующий тензору кривизны Римана. Две точки на геодезической $\gamma_{0}$ называются сопряженными, если существует ненулевое якобиево поле вдоль $\gamma_{0}$, обращающееся в нуль в обеих этих точках.



Т е о р е ма. Индекс квадратичной формы $\partial^{2} E$ в критич. точке $\gamma_{0}$ равен числу точек на геодезической $\gamma_{0}$, сопряженных вдоль $\gamma_{0}$ начальной точке $p=\gamma_{0}(0)$. При этом каждая сопряженная точка учитывается столько раз, какова ее кратность.



Поступая по аналогии с конечномерным случаем, рассмотрим функционал действия $E(\gamma)$ на пространстве $\Omega$. Этот функционал будет «функцией Морса», если все его критич. точки (то есть геодезические из точки $p$ в точку $q$ ) будут невырождены. Это будет в том и только в том случае, когда точки $p$ и $q$ не сопряжены друг другу (вдоль любой геодезической, соединяющей $p$ и $q$ ).



Т е о р е м а. Пусть $M-$ компактное многообразие, $p$ и $q-$ пара точек на $M$, не сопряженных ни вдоль какой геодезической плины, не превосходящей $a$. Тогда множество $\Omega^{a}=\left(E \leqslant a^{2}\right)$ гомотопически эквивалентно конечному клеточному комплексу, в к-ром каждая клетка размерности $\lambda$ взаимно однозначно соответствует геодезической индекса $\lambda$ (и длины, не превосходящей $a$ ).



Если теперь «устремить $a$ в бесконечность», то в результате получим, что пространство всех непрерывных путей на $M$, соединяющих $p$ с $q$, имеет гомотопич. тип счетного клеточного комплекса, в к-ром каждой геодезической из $p$ в $q$ с индексом $\lambda$ соответствует в точности одна клетка размерности $\lambda$.



Изложенная вариационная задача называется «3 а д а ч е й с закрепленными концами». Важное значение имеют «замкнутые экстремумы», возникающие в периодич. задаче В. и. в ц. Рассматривается пространство П $(M)$ всех замкнутых гладких кривых на многообразии $M$. Анализ периодич. задачи сложнее, чем изучение геодезических с закрепленными концами. Напр., задача о нахождении ко- личества простых (то есть некратных) замкнутых геодезических далеко не тривиальна (см. [1]). Каждая точка пространства $\Pi(M)$ есть гладкое отображение окружности $S^{1}$ в $M$. Множество линейно связных компонент пространства $\Pi(M)$ - это, по определению, множество «свободных» гомотопич. классов отображений $S^{1} \rightarrow M$.



Т е о р е м а. Если $M$ - компактное риманово гладкое многообразие отрицательной кривизны, то в каждом свободном гомотопич. классе замкнутых путей сушествует единственная замкнутая геодезическая.



Одномерная задача общего типа изучает величину (функционал)



\[

S(\gamma)=\int_{p}^{q} L(x(t), \dot{x}(t), t) d t

\]



и ее экстремали. Функция $L-$ лагранжиан. Экстремали $\gamma_{0}$ - это решения уравнений Эйлера-Лагранжа:



\[

\left(\partial L / \partial \dot{x}^{i}\right)^{\prime}=\partial L / \partial x^{i}

\]



Энергия - это $\dot{x}^{i} \partial L / \partial \dot{x}^{i}-L ;$ импульс - это $p_{i}=\partial L / \partial \dot{x}^{i}$; сила - это $f_{i}=\partial L / \partial x^{i}$.



Приме р 1. Траектории движения частицы массы $m$ в потенциальном силовом поле $f=-\operatorname{grad} U-$ это экстремали функционала



\[

S=\int\left(\frac{m}{z} \sum\left(\dot{x}^{i}\right)^{2}-U(x)\right) d x

\]



(принцип наименьшего действия).



П р и ме р 2. Движение релятивистской (свободной) частицы ненулевой массы $m>0$ определяется в пространстве Минковского с координатами $\left(x^{0}=c t, x^{1}, x^{2}, x^{3}\right)$ одним из двух функционалов (действий) на времениподобных кривых:



\[

S_{1}=\frac{m c}{z} \int\langle\dot{x}, \dot{x}\rangle d \tau

\]



где



\[

\begin{aligned}

\langle\dot{x}, \dot{x}\rangle & =\left(x^{0}\right)^{2}-\sum_{\alpha=1}^{3}\left(\dot{x}^{\alpha}\right)^{2} ; \\

S_{2} & =-m c \int\langle\dot{x}, \dot{x}\rangle^{1 / 2} d \tau .

\end{aligned}

\]



Мировые линии массивных частиц в пространстве Минковского - это экстремали функционалов $S_{1}$ и $S_{2}$ (релятивистский принцип наименьшего действия). Основным постулатом общей теории относительности является тезис (или гипотеза) Эйнштейна, согласно к-рому гравитационное поле - это метрика $g_{i j}$ в 4-мерном пространстве-времени [в каждой точке к-рого метрика имеет тип $(+---)$ ]. При отсутствии других сил пробная частица в гравитационном поле движется по времениподобным геодезическим, определяемым функционалом $S_{1}$. Частица нулевой массы движется по изотропным геодезическим.



Лагранжиан $L$ называется н е в ы ро жде н ны м, если $\operatorname{det}\left(\partial^{2} L / \partial \dot{x}^{i} \partial \dot{x}^{j}\right) \neq 0$, и называется с ильно невы рожд е н н ы М, если уравнения $p_{i}=\partial L(x, \dot{x}) / \partial \dot{x}^{i}$ можно взаимно однозначно и гладко решить в виде $\dot{x}^{i}=v^{i}(x, p)$ для всех $x, \dot{x}$. Г а м и л ь т о н и а но м $H(x, p)$ (для сильно невырожденных лагранжианов) называется энергия $E$, выраженная через $x$ и $p$.



Т е о ре м а. Пусть $L(x, \dot{x})$ - сильно невырожденный лагранжиан и $H(x, p)$ - гамильтониан, где $p=\partial L(x, \dot{x}) / \partial \dot{x}$ и $\dot{x}=v(x, p)$. Уравнения Эйлера-Лагранжа



\[

\partial L / \partial x-\frac{d}{d t}(\partial L / \partial \dot{x})=0

\]



эквивалентны уравнениям Гамильтона, в к-рых $x$ и $p$ считаются независимыми переменными:



\[

\dot{p}=-\partial H / \partial x, \quad \dot{x}=-\partial H / \partial p

\]





\end{document}